% ============================================================================
\section{Paradigm 2: Tsetlin Machines}
\label{sec:tsetlin}
% ============================================================================

% --- Intuition ---
\subsection{The Idea in Plain English}

Forget neurons entirely. Imagine instead a committee of \emph{rule
makers}. Each rule maker writes a simple if--then rule in plain
English:

\begin{quote}
\emph{``If pixel~(14,14) is dark AND pixel~(10,10) is bright AND
pixel~(20,5) is NOT bright, then it is probably a~3.''}
\end{quote}

Each rule maker starts with random rules and improves through
trial and error: rules that help classify correctly get reinforced;
rules that cause mistakes get weakened. No multiplication, no
gradients, no floating-point numbers. Just counting and logic.

This is the Tsetlin Machine~\cite{granmo2018tsetlin}.

\begin{analogy}
A traditional neural network is like an orchestra playing a
continuous melody---smooth, precise, nuanced. A Tsetlin Machine
is like a parliament voting on proposals. Each clause is a proposed
rule; the machine counts votes for and against each classification.
The learning process is the political negotiation: successful
proposals gain support; failed ones lose it.
\end{analogy}

% --- Mechanism ---
\subsection{How It Works}

\subsubsection{Step 1: Booleanise the Input}

Tsetlin Machines operate on Boolean (true/false) inputs. For MNIST,
each pixel is thresholded (bright/dark), producing 784 Boolean
features. For richer representation, each feature $x_k$ is
accompanied by its negation $\bar{x}_k$, doubling the input to
1568~literals.

\begin{figure}[htbp]
\centering
\begin{tikzpicture}[
  feat/.style={draw, minimum width=0.7cm, minimum height=0.6cm,
    font=\scriptsize\ttfamily, inner sep=0pt},
  lbl/.style={font=\small\sffamily}
]
  % Pixel value
  \node[lbl] at (-2, 0) {Pixel value:};
  \node[feat, fill=bitblue!30] at (0, 0) {183};

  % Arrow
  \draw[-{Stealth}, thick] (0.6, 0) -- (2, 0);
  \node[font=\scriptsize\sffamily, above] at (1.3, 0) {$> 127$?};

  % Boolean
  \node[feat, fill=bitgreen!25] (pos) at (3, 0.5) {$x_k = 1$};
  \node[feat, fill=bitorange!25] (neg) at (3, -0.5) {$\bar{x}_k = 0$};

  \draw[-{Stealth}] (2.1, 0.1) -- (pos.west);
  \draw[-{Stealth}] (2.1, -0.1) -- (neg.west);

  % Label
  \node[font=\scriptsize\sffamily, bitgray, right=0.3cm] at (neg.east)
    {Both $x_k$ and $\bar{x}_k$ are available as literals};
\end{tikzpicture}
\caption{Input booleanisation for Tsetlin Machines. Each pixel becomes
  two literals: the feature and its complement.}
\label{fig:tm-booleanise}
\end{figure}

\subsubsection{Step 2: Clauses---The ``Rules''}

A \concept{clause} is a conjunction (AND) of selected literals:

\begin{equation}
  C_j = x_3 \;\texttt{AND}\; \bar{x}_7 \;\texttt{AND}\; x_{42}
    \;\texttt{AND}\; \bar{x}_{100} \;\texttt{AND}\; \ldots
  \label{eq:clause}
\end{equation}

A clause outputs 1 only if \emph{all} its included literals are true.
Crucially, most literals are \emph{excluded}---a clause typically
selects only a small subset of the 1568 available literals.

\subsubsection{Step 3: Voting}

Clauses are assigned polarity: positive clauses vote \emph{for} a
class, negative clauses vote \emph{against}. The class score is:

\begin{equation}
  \text{score}(y) = \sum_{j \in \text{pos}(y)} C_j
    - \sum_{j \in \text{neg}(y)} C_j
  \label{eq:tm-voting}
\end{equation}

The predicted class is the one with the highest score.

\begin{figure}[htbp]
\centering
\begin{tikzpicture}[
  clause/.style={draw, rounded corners=3pt, minimum width=3.2cm,
    minimum height=0.7cm, font=\scriptsize\sffamily, align=center},
  vote/.style={-{Stealth}, thick},
  >=Stealth
]
  % Input
  \node[draw, rounded corners, fill=lightbg, minimum width=2cm,
    font=\small\sffamily] (input) at (0, 0) {Input (bits)};

  % Positive clauses
  \node[clause, fill=bitgreen!15] (c1) at (5, 1.5)
    {$C_1^+$: $x_3 \wedge \bar{x}_7 \wedge x_{42}$};
  \node[clause, fill=bitgreen!15] (c2) at (5, 0.5)
    {$C_2^+$: $x_{10} \wedge x_{15} \wedge \bar{x}_{200}$};

  % Negative clauses
  \node[clause, fill=bitorange!15] (c3) at (5, -0.5)
    {$C_3^-$: $x_1 \wedge x_{50} \wedge \bar{x}_{300}$};
  \node[clause, fill=bitorange!15] (c4) at (5, -1.5)
    {$C_4^-$: $\bar{x}_5 \wedge x_{77}$};

  % Connections
  \foreach \c in {c1,c2,c3,c4} {
    \draw[->] (input.east) -- (\c.west);
  }

  % Sum
  \node[draw, circle, minimum size=1cm, fill=bitblue!10,
    font=\small\sffamily\bfseries] (sum) at (9.5, 0) {$\Sigma$};

  \draw[vote, bitgreen] (c1.east) -- node[above, font=\scriptsize] {$+1$} (sum);
  \draw[vote, bitgreen] (c2.east) -- node[above, font=\scriptsize] {$+1$} (sum);
  \draw[vote, bitorange] (c3.east) -- node[below, font=\scriptsize] {$-1$} (sum);
  \draw[vote, bitorange] (c4.east) -- node[below, font=\scriptsize] {$-1$} (sum);

  % Output
  \node[font=\small\sffamily, right=0.6cm] at (sum.east)
    {Score for class ``3''};
\end{tikzpicture}
\caption{Tsetlin Machine clause voting. Each clause evaluates a
  conjunction of literals (an AND of selected input bits). Positive
  clauses vote for the class; negative clauses vote against. The net
  score determines the prediction.}
\label{fig:tm-voting}
\end{figure}

\subsubsection{Step 4: Learning with Tsetlin Automata}

This is where Tsetlin Machines are fundamentally different from neural
networks. There are no gradients. Instead, each literal in each clause
is controlled by a \concept{Tsetlin automaton}---a simple finite-state
machine with $2N$ states.

\begin{figure}[htbp]
\centering
\begin{tikzpicture}[
  state/.style={draw, circle, minimum size=0.6cm, font=\tiny\sffamily,
    inner sep=0pt},
  >=Stealth
]
  % States
  \foreach \i in {1,...,4} {
    \node[state, fill=bitorange!20] (ex\i) at (\i*1.2, 0)
      {$\i$};
  }
  \foreach \i in {5,...,8} {
    \node[state, fill=bitgreen!20] (in\i) at (\i*1.2, 0)
      {$\i$};
  }

  % Arrows between states
  \foreach \i [evaluate=\i as \j using int(\i+1)] in {1,...,7} {
    \pgfmathtruncatemacro{\ii}{\i}
    \pgfmathtruncatemacro{\jj}{\j}
    \ifnum\ii<4
      \draw[->] (ex\ii) -- (ex\jj);
    \else
      \ifnum\ii=4
        \draw[->] (ex4) -- (in5);
      \else
        \draw[->] (in\ii) -- (in\jj);
      \fi
    \fi
  }
  \foreach \i [evaluate=\i as \j using int(\i-1)] in {2,...,8} {
    \pgfmathtruncatemacro{\ii}{\i}
    \pgfmathtruncatemacro{\jj}{\j}
    \ifnum\ii>5
      \draw[->, bend right=30, dashed] (in\ii) to (in\jj);
    \else
      \ifnum\ii=5
        \draw[->, bend right=30, dashed] (in5) to (ex4);
      \else
        \draw[->, bend right=30, dashed] (ex\ii) to (ex\jj);
      \fi
    \fi
  }

  % Dividing line
  \draw[thick, bitgray, dashed] (5.4, -1) -- (5.4, 1);
  \node[font=\scriptsize\sffamily, bitorange] at (3, -1.3) {EXCLUDE literal};
  \node[font=\scriptsize\sffamily, bitgreen] at (7.8, -1.3) {INCLUDE literal};

  % Labels
  \node[font=\scriptsize\sffamily, bitgray] at (5.4, 1.3) {threshold};
\end{tikzpicture}
\caption{A Tsetlin automaton with $2N = 8$ states. States 1--4 mean
  ``exclude this literal from the clause''; states 5--8 mean ``include
  it.'' Learning nudges the state left or right based on feedback.}
\label{fig:tsetlin-automaton}
\end{figure}

The learning algorithm uses stochastic feedback:

\begin{description}[leftmargin=1em]
  \item[Type~I feedback] (reward correct clauses): Reinforce included
    literals that matched the input; randomly exclude literals that
    were not helpful. This \emph{specialises} clauses.
  \item[Type~II feedback] (penalise incorrect clauses): Include
    literals that would have falsified the clause (making it output~0
    instead of erroneously voting). This \emph{corrects} clauses.
\end{description}

\begin{engineernote}
The beauty of Tsetlin automata is that they are just
\emph{counters}. Increment or decrement a state by 1 based on a
coin flip and the classification result. No floating-point
arithmetic at all. The entire training loop is integer addition,
comparison, and random number generation.
\end{engineernote}

% --- Formalisation ---
\subsection{Mathematical Summary}
\label{sec:tm-math}

Let $\bm{x} \in \{0,1\}^n$ be a binary input vector. Define the
extended literal set $L = \{x_1, \bar{x}_1, x_2, \bar{x}_2, \ldots,
x_n, \bar{x}_n\}$ with $|L| = 2n$.

Each clause $C_j$ is defined by a subset $S_j \subseteq L$ and
computes:
\begin{equation}
  C_j(\bm{x}) = \bigwedge_{\ell \in S_j} \ell
  \label{eq:tm-clause-formal}
\end{equation}

For a $K$-class problem with $m$ clauses per class (half positive,
half negative), the output for class $k$ is:
\begin{equation}
  f_k(\bm{x}) = \sum_{j=1}^{m/2} C_j^{+,k}(\bm{x})
    - \sum_{j=1}^{m/2} C_j^{-,k}(\bm{x})
  \label{eq:tm-class-score}
\end{equation}

The predicted class is $\hat{y} = \arg\max_k f_k(\bm{x})$.

The threshold parameter $T$ controls clause voting saturation (scores
are clipped to $[-T, T]$), and the specificity parameter $s$ controls
how aggressively clauses specialise during Type~I feedback.

\subsection{Performance on MNIST}
\label{sec:tm-mnist}

Tsetlin Machines achieve remarkable results on MNIST~\cite{lecun1998mnist}:

\begin{table}[htbp]
\centering
\caption{Tsetlin Machine results on MNIST. Results from the original
  TM paper~\cite{granmo2018tsetlin} and the Convolutional TM
  paper~\cite{granmo2019convtm}.}
\label{tab:tm-mnist}
\begin{tabular}{@{}lrrl@{}}
  \toprule
  \textbf{Configuration}     & \textbf{Accuracy} & \textbf{Clauses} & \textbf{Source} \\
  \midrule
  Basic TM (4{,}000/class)   & 98.2\%            & 40{,}000 & \cite{granmo2018tsetlin} \\
  Convolutional TM            & 99.3\%            & 80{,}000 & \cite{granmo2019convtm} \\
  \bottomrule
\end{tabular}
\end{table}

\subsection{Strengths and Limitations}

\begin{description}[leftmargin=1em]
  \item[Strengths:]
    No floating-point arithmetic at all (training or inference),
    fully interpretable (you can read the learned rules),
    convergence proofs exist~\cite{granmo2018tsetlin},
    naturally parallel (clauses are independent),
    robust to noise.

  \item[Limitations:]
    Requires many clauses for complex tasks (millions for high
    accuracy), limited to propositional logic (no hierarchical
    feature extraction by default), hyperparameter sensitivity
    ($T$ and $s$), less mature ecosystem than neural networks.

  \item[Open question:]
    In the standard formulation, the number of clauses and the literal
    pool are both fixed before training begins. The automata learn
    \emph{which} literals to include in each clause, but cannot create
    new clauses at runtime. Clause weighting variants can suppress
    unhelpful clauses (effectively ``turning them off''), and drop-clause
    regularisation randomly disables clauses during training (analogous
    to dropout in neural networks). However, dynamically \emph{growing}
    the clause set during training---adding clauses when the model
    detects that its current capacity is insufficient---remains an
    unexplored research direction.
\end{description}
